\documentclass[11pt,a4paper]{article}

%%%%%%%%%%%%%%%%%%%%
% STANDARD PACKAGES
%%%%%%%%%%%%%%%%%%%%
\usepackage[utf8]{inputenc}
\usepackage[T1]{fontenc}
\usepackage{lmodern}
\usepackage[margin=1in]{geometry}
\usepackage{amsmath,amssymb,amsfonts,bm}
\usepackage{siunitx}
\usepackage{physics}
\usepackage{graphicx}
\usepackage{booktabs}
\usepackage[numbers,super,sort&compress]{natbib}
\usepackage[colorlinks=true,allcolors=blue]{hyperref}

% Fix for siunitx and physics package conflict
\AtBeginDocument{\RenewCommandCopy\qty\SI}

% Covariant d'Alembertian (unused but kept)
\newcommand{\wave}{\nabla^\mu\nabla_\mu}

%%%%%%%%%%%%%%%%%%%%
% TITLE & AUTHOR
%%%%%%%%%%%%%%%%%%%%
\title{\textbf{A Mechanical Route to \texorpdfstring{$G$}{G}: Matching Field Stress--Energy to the Einstein--Hilbert Coupling}}

\author{
    Robin Victor Skavberg\thanks{ORCID: \href{https://orcid.org/0009-0003-9126-6196}{0009-0003-9126-6196}} \\
    \small Independent Researcher \\
    \small Calgary, Alberta, Canada
}

\date{January 26, 2026}

%%%%%%%%%%%%%%%%%%%%
% BEGIN DOCUMENT
%%%%%%%%%%%%%%%%%%%%
\begin{document}

\maketitle

\begin{abstract}
We develop a mechanical route to Newton’s gravitational coupling $G$ by relating the macroscopic field stress--energy to the geometric stiffness of the Einstein--Hilbert (EH) action. In the long--wavelength regime of a homogeneous, barotropic background, the bulk modulus satisfies $K=\rho_m c_s^2$; identifying the signal speed with $c$ yields $K=\rho_m c^2$. We then show that an EH term $\int\!\sqrt{-g}\,R$ arises with a calculable coefficient in the one--loop effective action, so that matching the geometric stiffness $c^4/(8\pi G)$ to the mechanical stiffness gives $G=c^2/(8\pi\rho_m)$. Independently, we verify the observational identity $G=\Lambda c^4/(8\pi u_\nu)$, where $\Lambda$ is the cosmological constant and $u_\nu$ is the corresponding vacuum stress--energy density, providing a parameter--free cross--check.

To connect with microphysics and phenomenology, we outline three realizations: (A) an elastic curvature response consistent with thermodynamic and induced--gravity perspectives; (B) a superfluid dark--matter realization in which phonon excitations couple to baryons and exert a literal force; and (C) a single--field EFT of dark energy with explicit matter couplings that generate an effective force and $G_{\rm eff}(t,k)$. Our baseline adopts (A) (no extra force) and recovers $\Lambda$CDM at linear order; (B) and (C) are presented as optional model ingredients with distinct, testable signatures.
\end{abstract}

\section{Introduction}
Analogue--gravity studies show that hydrodynamic media can reproduce relativistic kinematics via an effective metric,\cite{Unruh1981,Visser1998,Barcelo2011} suggesting a conservative route to relate gravitational dynamics to mechanical properties. In this work, we develop a mechanical route to $G$ by modeling the background as a barotropic medium. Identifying the signal speed with $c$ fixes the field pressure to $K=\rho_m c^2$, which we match to the geometric stiffness multiplying curvature in the Einstein--Hilbert (EH) action.\cite{Einstein1916,Hilbert1915}

\section{Assumptions and Regime of Validity}
We assume: (i) long--wavelength hydrodynamics ($k\xi\ll1$); (ii) homogeneous, barotropic background with mass density $\rho_m$; (iii) phase coherence; (iv) identification of acoustic and electromagnetic null cones in the IR ($c_s=c$).

\section{Field Pressure: Three Realizations and Couplings}
\label{sec:three-realizations}
We use ``field pressure'' as shorthand for the macroscopic stress--energy that sets a response scale for spacetime. The way this response acts on matter depends on microphysics and couplings. We distinguish three realizations:

\paragraph{(A) Elastic curvature response (baseline).}
With minimal coupling and a cosmological constant background, the field pressure fixes the medium's elastic response; matter follows geodesics of the curved metric. The ``push back'' is spacetime curvature---an equation-of-state relation---not a direct pressure force on baryons.\cite{Jacobson1995,Padmanabhan2004,Sakharov1967}

\paragraph{(B) Superfluid dark matter with phonon force (optional).}
If the field is a dark-matter superfluid in galaxies, the phonon mode $\theta$ can couple to baryons, e.g.
$\mathcal{L}_{\rm int}=\alpha\,\theta\,\rho_b$, so baryons feel an additional force
$\mathbf{a}_{\rm phonon}=-\alpha\,\nabla\theta$ on top of geodesic motion. Here pressure/stress gradients transmit a literal force in the superfluid phase, while cosmology remains CDM-like.\cite{BerezhianiKhoury2015,KhouryPLB2015}

\paragraph{(C) Single-field EFT of dark energy with explicit couplings (optional).}
In a general single-field EFT, the background field supplies the pressure/energy density, and matter may couple conformally/disformally to the scalar sector. At the perturbative level this produces an effective force on baryons and a scale- and time-dependent $G_{\rm eff}(t,k)$, subject to stability and observational constraints; setting extra quadratic operators to zero reproduces $\Lambda$CDM at linear order.\cite{Gubitosi2013}

\medskip
\noindent\textbf{Constraint summary.} Any direct baryon--vacuum interaction producing a fifth force is excluded at the level $|\alpha|\!\lesssim\!10^{-13}$ (and below) by precision equivalence–principle tests unless it is screened or emergent; therefore, in the baseline scenario baryons couple to vacuum stress--energy only through spacetime curvature, while literal forces arise only in the optional superfluid or EFT realizations.\cite{EotWashEP,AsteroidFifthForce,Jacobson1995,Gubitosi2013,BerezhianiKhoury2015}

\section{Hydrodynamic Framework}
At the microscopic level, we take a complex scalar condensate described by the Gross--Pitaevskii (GP) equation:\cite{Gross1961,Pitaevskii1961}
\begin{equation}
i\hbar\,\partial_t\Psi=\left[-\frac{\hbar^2}{2m}\nabla^2 + V_{\rm ext} + g|\Psi|^2\right]\Psi.
\end{equation}
Writing $\Psi=\sqrt{n}\,e^{i\theta}$ gives hydrodynamics for number density $n$ and velocity $\vb{v}=(\hbar/m)\nabla\theta$. Linear perturbations propagate on an acoustic metric $ds^2\propto -c_s^2\,dt^2 + d\vb{x}^2$.

\section{Mechanical Stiffness and Geometric Coupling}

\subsection{Bulk Modulus from the Equation of State}
For the GP fluid, the pressure is $P=\tfrac12 g n^2$. We define $\rho_m \equiv mn$ as the mass density of the condensate (where $m$ is the constituent mass and $n$ is the number density). The sound speed is determined by the derivative of pressure with respect to the mass density:\cite{LandauLifshitz}
\begin{equation}
c_s^2 = \frac{dP}{d\rho_m} = \frac{d(\tfrac{1}{2}gn^2)}{d(mn)} = \frac{gn}{m}.
\end{equation}
The bulk modulus (mechanical stiffness) is $\beta \equiv \rho_m \frac{dP}{d\rho_m}$. Setting $c_s=c$ yields:
\begin{equation}
K \equiv \beta = \rho_m c^2.
\end{equation}

\subsection{Matching to Einstein’s Coupling}
The EH action has the prefactor $c^4/(16\pi G)$, corresponding to a geometric stiffness $c^4/(8\pi G)$. 
Equating mechanical and geometric stiffnesses gives:
\begin{equation}
\label{eq:Gfromrho}
\rho_m c^2 = \frac{c^4}{8\pi G}
\;\;\implies\;\;
G = \frac{c^2}{8\pi \rho_m}.
\end{equation}

\subsection{Verification and Significance}
The identity
\begin{equation}
G=\frac{\Lambda c^4}{8\pi u_\nu}
\end{equation}
follows from the standard equivalence between the cosmological constant term
$\Lambda g_{\mu\nu}$ and a Lorentz-invariant vacuum stress--energy
$T^{(\Lambda)}_{\mu\nu}=-u_\nu g_{\mu\nu}$ with $u_\nu=\rho_\Lambda c^2=\Lambda c^4/(8\pi G)$.
Using observational values for $\Lambda$ and $u_\nu$ yields
$G\simeq \SI{6.67e-11}{m^3.kg^{-1}.s^{-2}}$, a parameter-free cross-check.
In our constitutive interpretation, the background vacuum stress--energy (associated with $\Lambda$)
sets the elastic response scale of spacetime; $G$ is the corresponding modulus.
Realizations (B) and (C) introduce explicit matter couplings through which pressure/stress gradients
can exert a literal force on baryons; our baseline analysis adopts (A) with no extra force and hence
reduces to $\Lambda$CDM at linear order.\cite{Jacobson1995,Gubitosi2013,BerezhianiKhoury2015}

\paragraph{Interpretive consistency.}
Our viewpoint is conceptually consistent with: (i) Jacobson’s reading of the Einstein equations
as an \emph{equation of state}; (ii) Sakharov’s induced-gravity picture in which the
Einstein--Hilbert term encodes a \emph{metric elasticity} generated by vacuum fluctuations; and
(iii) Padmanabhan’s paradigm of \emph{gravity as spacetime elasticity}.
Accordingly, we take the background vacuum stress--energy to set the \emph{macroscopic stiffness}
of the gravitational medium. Newton’s constant $G$ then appears as a \emph{constitutive parameter}
(effective modulus) quantifying spacetime’s resistance to curvature; matter curves spacetime, and
the elastic response is fixed by this modulus.\cite{Jacobson1995,Sakharov1967,Padmanabhan2004}

\subsection{Novel Correlation and Derivation}
While the Gross--Pitaevskii equation describes a condensed medium and the heat-kernel expansion organizes the gravitational effective action, these frameworks are usually treated separately. Our contribution is the \emph{stiffness matching} condition: in the IR (long-wavelength) regime, the mechanical bulk modulus is matched to the EH geometric stiffness, turning $G$ from an empirical constant into a constitutive parameter fixed by the background state. This matching is compatible with an elastic/thermodynamic reading of spacetime and admits microphysical realizations with and without explicit forces on baryons.

\section{Dimensional Consistency and Physical Interpretation}
To ensure the mathematical validity of the proposed mechanical route, we perform a dimensional analysis of the identity $G = \Lambda c^4 / (8\pi u_\nu)$. As shown in Table~\ref{tab:dimensions}, the units of the derived coupling perfectly match the standard dimensions of Newton's gravitational constant.

\begin{table}[h]
\centering
\caption{Dimensional Analysis of the Mechanical Route to $G$}
\label{tab:dimensions}
\vspace{2mm}
\begin{tabular}{@{}llll@{}}
\toprule
\textbf{Quantity} & \textbf{Symbol} & \textbf{SI Units} & \textbf{Dimensions} \\ \midrule
Gravitational Constant & $G$ & \si{m^3.kg^{-1}.s^{-2}} & $[L]^3 [M]^{-1} [T]^{-2}$ \\
Speed of Light & $c$ & \si{m.s^{-1}} & $[L] [T]^{-1}$ \\
Vacuum Energy Density & $u_\nu$ & \si{kg.m^{-1}.s^{-2}} (Pa) & $[M] [L]^{-1} [T]^{-2}$ \\
Cosmological Constant & $\Lambda$ & \si{m^{-2}} & $[L]^{-2}$ \\
Condensate Density & $\rho_m$ & \si{kg.m^{-3}} & $[M] [L]^{-3}$ \\ \bottomrule
\end{tabular}
\end{table}

\begin{table}[h]
\centering
\caption{Algebraic Breakdown of Dimensional Consistency}
\label{tab:algebraic_dims}
\vspace{2mm}
\begin{tabular}{@{}ll@{}}
\toprule
\textbf{Component} & \textbf{Dimensional Operation} \\ \midrule
Numerator ($[\Lambda] \times [c]^4$) & $[L]^{-2} \times ([L][T]^{-1})^4 = [L]^2 [T]^{-4}$ \\[3pt]
Denominator ($[u_\nu]$) & $[M] [L]^{-1} [T]^{-2}$ \\[3pt]
\textbf{Final Ratio} ($\frac{\text{Numerator}}{\text{Denominator}}$) & $\frac{[L]^2 [T]^{-4}}{[M] [L]^{-1} [T]^{-2}} = \mathbf{[L]^3 [M]^{-1} [T]^{-2}}$ \\ \midrule
\textbf{Identity Check} & $\equiv [G]$ \\ \bottomrule
\end{tabular}
\end{table}

\subsection*{Physical Significance of the Result}
The dimensional alignment supports the following physical interpretations:
\begin{itemize}
    \item \textbf{Energy density as field pressure:} The vacuum energy density $u_\nu$ (Pa) sets the macroscopic scale of the medium’s resistance; in the baseline (A) it fixes the elastic curvature response, while in (B) and (C) pressure/stress gradients can also exert literal forces when couplings are present.
    \item \textbf{Curvature as deformation:} The cosmological constant $\Lambda$ (inverse length squared) characterizes a uniform curvature scale of spacetime.
    \item \textbf{Gravity as an emergent stiffness:} The gravitational constant $G$ is the effective modulus linking stress--energy to curvature in the IR description.
\end{itemize}

\section{Phenomenological Outlook (Optional Realizations)}
The superfluid realization (B) predicts halo--scale signatures (e.g., phonon--mediated forces, cores, and phase--dependent dynamics), while the EFT realization (C) organizes linear--regime deviations via $G_{\rm eff}(t,k)$ under stability and observational constraints. Our baseline (A) introduces no extra force and reproduces standard $\Lambda$CDM growth and lensing at linear order. A dedicated exploration of (B)/(C) against data is left for future work.

\section{Conclusion}
We have shown that Newton's gravitational constant $G$ can be interpreted as a constitutive parameter by matching the mechanical bulk modulus of a barotropic medium to the geometric stiffness of the Einstein--Hilbert action. An observational identity using the cosmological constant and vacuum stress--energy density provides a parameter--free cross--check of the numerical value. The framework is consistent with an elastic/thermodynamic reading of spacetime and admits realizations with and without explicit forces on baryons. Our baseline recovers $\Lambda$CDM at linear order; optional superfluid and EFT realizations offer distinct, testable signatures.

\section*{Declarations}
\subsection*{AI Transparency Statement}
The fundamental theoretical framework, the correlation matching mechanical bulk modulus to geometric stiffness, and the conceptual interpretation of $G$ are the original work of the author. Generative AI tools were used for LaTeX formatting, bibliography organization, and prose polishing.

\appendix

\section{One--loop derivation of the EH term}
\label{app:one-loop}
We integrate out field fluctuations. The divergent part of the one-loop effective action $\Gamma_{1}$ is given by the Schwinger proper-time representation:\cite{Vassilevich03,Schwinger1951}
\begin{equation}
\Gamma_{1} = -\frac{1}{2(4\pi)^2} \int d^4x \sqrt{-g} \int_{1/\Lambda_{UV}^2}^\infty \frac{ds}{s^3} e^{-sM_{\rm eff}^2} [a_0 + a_1 s + a_2 s^2].
\end{equation}
The term $a_1 = (\tfrac{1}{6}-\xi)R$ induces the EH action:
\begin{equation}
S_{\rm ind} \supset \frac{(\tfrac{1}{6}-\xi)\Lambda_{UV}^2}{32\pi^2} \int d^4x \sqrt{-g} R.
\end{equation}
Matching this to $S_{\rm EH} = \frac{c^4}{16\pi G} \int \sqrt{-g} R$ anchors $G$ to the medium's cutoff $\Lambda_{UV}$ in the induced-gravity picture.

\section{Stress Test: Objections and Responses (FAQ)}
\label{app:faq}
\textbf{Q1: Is this just GR with $\Lambda$? Where is the novelty?}\\
\emph{A:} The field equations are standard. The novelty is interpretive and structural:
we treat the background vacuum stress--energy as primary and read $G$ as a macroscopic modulus,
aligned with thermodynamic/induced--gravity perspectives.\cite{Jacobson1995,Sakharov1967,Padmanabhan2004}

\textbf{Q2: Does ``field pressure'' push directly on matter?}\\
\emph{A:} In the baseline (A), no: the response is elastic curvature (geodesic motion).
In realizations (B) and (C), explicit couplings let pressure/stress gradients exert a literal force
(phonon-mediated in a superfluid; scalar--matter couplings in EFT).\cite{BerezhianiKhoury2015,Gubitosi2013}

\textbf{Q3: Does this imply a varying $G$?}\\
\emph{A:} Not in the baseline, which assumes a homogeneous, quasi-constant background ($w\simeq -1$).
Generalizations with $G_{\rm eff}(t,k)$ fit naturally into EFT-of-DE and are constrained by growth/lensing;
we leave quantitative fits to future work.\cite{Gubitosi2013}

\section{Illustrative Extension: Evolution Across Epochs (Hypothetical)}
\label{app:evolution}
The mechanical interpretation suggests a possible relation between $G$ and the medium density $\rho_m$. If $\rho_m$ were to track cosmic evolution, one could obtain a varying $G\propto \rho_m^{-1}$. This is a \emph{hypothetical} extension subject to stringent constraints (BBN, CMB, local tests); we do not claim such variations are realized in nature. The table below illustrates the scaling logic only.

\begin{table}[h]
\centering
\caption{Illustrative scaling in a hypothetical tracking regime}
\label{tab:epochs}
\vspace{2mm}
\begin{tabular}{@{}llll@{}}
\toprule
\textbf{Epoch} & \textbf{Redshift ($z$)} & \textbf{Rel. Density $\rho_m / \rho_0$} & \textbf{Rel. Coupling $G/G_0$} \\ \midrule
Radiation Dom. & $> 3000$ & $\sim 10^{10}$ & $\sim 10^{-10}$ \\
Matter Dom. & $\sim 1$ & $8$ & $0.125$ \\
Present Day & $0$ & $1$ & $1$ \\ \bottomrule
\end{tabular}
\end{table}

%%%%%%%%%%%%%%%%%%%%
% REFERENCES
%%%%%%%%%%%%%%%%%%%%
\begin{thebibliography}{99}

\bibitem{Unruh1981} W. G. Unruh, Phys. Rev. Lett. \textbf{46}, 1351 (1981).
\bibitem{Visser1998} M. Visser, Class. Quantum Grav. \textbf{15}, 1767 (1998).
\bibitem{Barcelo2011} C. Barcel\'{o}, S. Liberati, and M. Visser, Living Rev. Relativ. \textbf{14}, 3 (2011).
\bibitem{Einstein1916} A. Einstein, Annalen der Physik \textbf{354}, 769 (1916).
\bibitem{Hilbert1915} D. Hilbert, Nachr. Ges. Wiss. G\"ottingen, Math.-Phys. Kl. 395 (1915).
\bibitem{Gross1961} E. P. Gross, Nuovo Cimento \textbf{20}, 454 (1961).
\bibitem{Pitaevskii1961} L. P. Pitaevskii, Sov. Phys. JETP \textbf{13}, 451 (1961).
\bibitem{LandauLifshitz} L. D. Landau and E. M. Lifshitz, \textit{Fluid Mechanics}, 2nd ed. (Pergamon, London, 1987).
\bibitem{Vassilevich03} D. V. Vassilevich, Phys. Rep. \textbf{388}, 279 (2003).
\bibitem{Schwinger1951} J. Schwinger, Phys. Rev. \textbf{82}, 664 (1951).

% Cosmology parameters (optional; kept for completeness)
\bibitem{PDG2023} O. Lahav and A. R. Liddle, in \textit{Review of Particle Physics}, PTEP \textbf{2022}, 083C01 (2022).
\bibitem{BoylanKolchin2023} M. Boylan-Kolchin, Nature Astron. \textbf{7}, 731 (2023).

% Three realizations anchors:
\bibitem{Jacobson1995} T. Jacobson, Phys. Rev. Lett. \textbf{75}, 1260 (1995). % Einstein eq. as equation of state
\bibitem{Sakharov1967} A. D. Sakharov, Sov. Phys. Dokl. \textbf{12}, 1040 (1968). % Induced gravity
\bibitem{Padmanabhan2004} T. Padmanabhan, Int. J. Mod. Phys. D \textbf{13}, 2293 (2004). % Gravity as elasticity

\bibitem{BerezhianiKhoury2015} L. Berezhiani and J. Khoury, Phys. Rev. D \textbf{92}, 103510 (2015). % Superfluid DM
\bibitem{KhouryPLB2015} L. Berezhiani and J. Khoury, Phys. Lett. B \textbf{753}, 639 (2016). % Phenomenology

\bibitem{Gubitosi2013} G. Gubitosi, F. Piazza, and F. Vernizzi, JCAP \textbf{02}, 032 (2013). % EFT of DE

% Fifth-force/EP constraints cited in the constraint summary:
\bibitem{EotWashEP} E\"ot-Wash Group, ``Equivalence Principle,'' University of Washington (accessed 2026). Available online.
\bibitem{AsteroidFifthForce} Y.-D. Tsai, Y. Wu, S. Vagnozzi, and L. Visinelli,
JCAP \textbf{04}, 031 (2023).

\end{thebibliography}

\end{document}